% =====================================================================
%  Chapter 23 opener -- "Chapter at a Glance", three reading paths + the
%  ladder staircase, THE running example (the shared 16-slot array
%  A=(3,2,4,5,1,1,5,3,7,0,2,4,1,6,3,2), sum 49 -- this chapter is array/
%  function based, so beats traces run on slices of A and the Li Chao
%  sections trace on an explicit x-grid / line set), the three engines
%  that rescue the non-composable updates Chapters 19-22 handed off, the
%  who-owns-what comparison, the recognition guide, a constraint->engine
%  table, and notation.
%  Exports fig:c29-gallery, fig:c29-taxonomy.
%  No \term{segment tree beats} here -- that term is indexed by Ch.20; use
%  \emph{beats} + a prose cite to avoid the document-wide INDEX overflow.
% =====================================================================

\begin{chaptertheorybox}{Chapter at a Glance}
Chapter~19 folded one aggregate over a range and edited a point.  Chapter~20
rewrote a whole range by pushing a \emph{single} lazy tag down, as long as the
updates form a \emph{composable monoid}.  Chapter~22 tiled a tree path into
\(O(\log n)\) chain ranges and reused that tag.  Every one of them rests on the
same promise: a range write is an operation you can \emph{compose} and
\emph{defer}.  Break the promise and all three fall silent.  \emph{Set every
element in \([l,r]\) to \(\min(a_i,x)\)} does not commute with \emph{add}: the
tag you would defer depends on the values it lands on.  \emph{Reduce every
element modulo \(x\)} shrinks only the elements above \(x\).  \emph{Take a point
maximum over a moving family of lines} has no fixed array at all.  This chapter
is the gallery of the three engines that still answer those questions in
almost-logarithmic time.

\textbf{You will learn to}
\begin{itemize}[nosep]
  \item build a \term{Li Chao line container} -- insert a line and query the
        pointwise extremum over a fixed domain in \(\bigO{\log C}\), the descent
        that keeps the mid-winner and recurses the loser to one side
        (\(\secref{c29-lichao}\)); \hfill\lv{3}--\lv{4}
  \item ride that container \emph{down a tree} with insert-on-enter and
        rollback-on-exit, so each node minimises over its whole ancestor path
        (\(\secref{c29-path}\), SPOJ~COT6); \hfill\lv{4}--\lv{5}
  \item run \emph{segment-tree beats}: store \((\max,\mathrm{smax},\mathrm{cnt})\) per
        node and a range \(\mathit{chmin}\) that recurses only when it must, plus
        the amortised blending proof; and the random-data \term{Chtholly tree}
        (\term{ODT}) sibling (\(\secref{c29-beats}\), CF~896C); \hfill\lv{4}--\lv{5}
  \item let a segment-tree node store a \emph{transform} (an affine map or a
        matrix) and merge by \term{function composition}, folding a product of
        maps over a range (\(\secref{c29-func}\), CF~750E); \hfill\lv{5}
  \item read a specification and pick the \emph{right} advanced tree -- or hand
        it back to Ch.~19/20/21/22 -- off the ownership map
        (\(\secref{c29-boundary}\), the pick table). \hfill\lv{5}
\end{itemize}

\textbf{Prerequisites.} Chapter~19 (\chref{c13}) -- the segment tree whose
build/merge skeleton every engine here reuses.  Chapter~20 (\chref{c15}) -- the
lazy tag as a \emph{composition monoid}; \(\secref{c29-func}\) is its general
case and \(\secref{c29-beats}\) is what you do when the update is \emph{not}
composable.  Chapter~22 (\chref{c16}) -- the reason the whole ancestor path can
be queried at once (\(\secref{c29-path}\) contrasts it with chain tiling).
\end{chaptertheorybox}

% ---------------- Reading paths + ladder position ----------------
\begin{chaptertheorybox}{How to Read This Chapter (three paths)}
Level badges \(\lv3\)--\(\lv6\): \(\lv3\) is inserting a couple of lines into a
Li Chao container by hand and reading a point query off it; \(\lv4\) is the
insert-descent written correctly, or a single range \(\mathit{chmin}\) that
prunes; \(\lv{5}\) is the amortised beats proof, or a container that rolls back
along a DFS; \(\lv{6}\) is merging two Li Chao trees small-to-large on a tree DP
(the Summit).

\textbf{Line-container path (\(\lv3\)--\(\lv{5}\), ``queries over a family of
lines'').} \secref{c29-lichao} in full -- the descent, the \(\bigO{\log C}\)
bound, CF~631E -- then \secref{c29-path}, where the same container rides a DFS
with rollback (COT6), and finally the Summit \(\probref{c29-summit}\) where the
containers \emph{merge}.

\textbf{Amortised-updates path (\(\lv4\)--\(\lv{5}\), ``the write is not a
tag'').} \secref{c29-beats}: the \((\max,\mathrm{smax},\mathrm{cnt})\) node, the range
\(\mathit{chmin}\) with its blending proof, the modulo prune (with CF~438D in
prose, owned by Ch.~19), and the Chtholly tree that trades worst-case guarantees
for random-data amortisation (CF~896C).

\textbf{Composition path (\(\lv{4}\)--\(\lv{6}\), ``a node is a function'').}
\secref{c29-func} -- affine and matrix transforms composed on merge, back-citing
Chapter~20's monoid, capped by the \((\min,+)\)-matrix automaton tree of CF~750E.
Then \secref{c29-boundary}: the pick table that decides \emph{which} engine, and
hands dynamic hulls to Ch.~25 (\chref{c23}) and offline to Ch.~27 (\chref{c21}).
\end{chaptertheorybox}

\begin{figure}[H]
\centering
\begin{tikzpicture}[diagram, x=9mm, y=7mm]
  \foreach \st/\lab in {1/lv1,2/lv2,3/lv3,4/lv4,5/lv5,6/lv6} {
    \fill[\lab!22, draw=\lab, line width=0.8pt] (\st-1,0) rectangle (\st,\st*0.52);
    \node[font=\scriptsize\bfseries, text=\lab] at (\st-0.5,\st*0.52-0.42) {L\st};}
  \node[annot, anchor=south west, font=\scriptsize]
    at (0.05,0.62) {23.1 insert two lines; read a point query by hand};
  \node[annot, anchor=south west, font=\scriptsize]
    at (1.05,1.14) {23.1 the insert-descent; O(log C); solved CF 631E};
  \node[annot, anchor=south west, font=\scriptsize]
    at (2.05,1.66) {23.3 range chmin prunes on the max; the (max,smax,cnt) node};
  \node[annot, anchor=south west, font=\scriptsize]
    at (3.05,2.18) {23.2 container on a DFS with rollback; solved SPOJ COT6};
  \node[annot, anchor=south west, font=\scriptsize, text=lv5]
    at (4.05,2.70) {23.3 beats amortised proof + Chtholly; 23.4 compose-on-merge};
  \node[annot, anchor=south west, font=\scriptsize, text=lv6]
    at (5.05,3.22) {23.5 the pick table -- which engine; Summit CF 932F};
\end{tikzpicture}
\caption*{Chapter~23 on the six-step ladder: reading a point query off a couple
of inserted lines is \(\lv3\)--\(\lv{4}\); a correct insert-descent or a pruning
range \(\mathit{chmin}\) is \(\lv{4}\); the rollback-on-DFS container, the beats
amortised proof, and matrix composition are \(\lv{5}\); merging Li Chao trees on
a tree DP (the Summit) is step~6.}
\end{figure}

% ---------------- The running example ----------------
\subsection*{The Running Example: An Array, and the Updates That Break a Tag}

Chapters~18--20 share the running array
\[
  A=(3,2,4,5,1,1,5,3,7,0,2,4,1,6,3,2),\qquad \textstyle\sum A = 49,
\]
over indices \(1..16\).  This chapter is array/function based, so \emph{every}
hand trace uses \(A\) (or a small \(x\)-grid of lines), never the tree \(T\) of
Chapter~22.  Three updates make the point of the whole chapter:

\begin{itemize}[nosep, leftmargin=1.6em]
  \item \(\mathit{chmin}(l,r,x)\): replace each \(a_i\) in \([l,r]\) by
        \(\min(a_i,x)\).  On \(A[1..8]=(3,2,4,5,1,1,5,3)\) with \(x=3\) this
        gives \((3,2,3,3,1,1,3,3)\) -- the two \(5\)s and the \(4\) shrink, the
        rest are untouched.  Not composable with a later add, so Ch.~20's tag
        cannot own it: this is \emph{beats} (\(\secref{c29-beats}\)).
  \item \(a_i \mathbin{\%} x\): reduce in place.  A \(5\%3=2\) but a \(1\%3=1\)
        -- only elements \(\ge x\) move, and \(\ge 2x\) shrink by \(\ge\) half,
        the amortised handle (\(\secref{c29-beats}\), CF~438D in prose).
  \item \(\max_j (m_j x + b_j)\) at a query point \(x\): no array at all, a
        family of lines whose upper envelope you probe -- the Li Chao container
        (\(\secref{c29-lichao}\)).
\end{itemize}

\begin{figure}[H]
\centering
\begin{tikzpicture}[diagram, scale=0.62, transform shape,
   cell/.style={rectangle, draw, minimum width=7mm, minimum height=7mm,
                font=\scriptsize\ttfamily, inner sep=0pt},
   cut/.style={rectangle, draw=warningred, fill=warningred!12, minimum width=7mm,
                minimum height=7mm, font=\scriptsize\ttfamily, inner sep=0pt}]
  \node[font=\scriptsize\bfseries, anchor=east] at (0,1.1) {\(A\)};
  \foreach \val [count=\k] in {3,2,4,5,1,1,5,3,7,0,2,4,1,6,3,2} {
     \ifnum\k<9 \node[cell] (a\k) at (\k-0.5,1.1) {\(\val\)};
     \else \node[cell] (a\k) at (\k-0.5,1.1) {\(\val\)}; \fi
  }
  \node[font=\scriptsize, text=darkgray, anchor=east] at (0,0.2) {\(\mathit{chmin}\)};
  \foreach \val/\k in {3/1,2/2,3/3,3/4,1/5,1/6,3/7,3/8} {
     \ifnum\k=3 \node[cut] at (\k-0.5,0.2) {\(\val\)};
     \else\ifnum\k=4 \node[cut] at (\k-0.5,0.2) {\(\val\)};
     \else\ifnum\k=7 \node[cut] at (\k-0.5,0.2) {\(\val\)};
     \else \node[cell] at (\k-0.5,0.2) {\(\val\)}; \fi\fi\fi
  }
  \node[font=\scriptsize, anchor=west] at (8.6,0.2) {\(A[9..16]\) untouched};
  \draw[warningred,->,line width=0.9pt] (a3.south) -- (3-0.5,0.55);
  \draw[warningred,->,line width=0.9pt] (a4.south) -- (4-0.5,0.55);
  \draw[warningred,->,line width=0.9pt] (a7.south) -- (7-0.5,0.55);
\end{tikzpicture}
\caption{The update that breaks a lazy tag.  Top: \(A[1..8]=(3,2,4,5,1,1,5,3)\).
Bottom: after \(\mathit{chmin}(1,8,3)\) the entries \(4,5,5\) clamp to \(3\)
(red) and the rest stay.  A single tag cannot express ``rewrite only the ones
above \(3\)'' -- the change depends on each value -- which is exactly why this
family of updates needs the beats bookkeeping of \(\secref{c29-beats}\).}
\label{fig:c29-gallery}
\end{figure}

% ---------------- Which engine, cost chart ----------------
\subsection*{Which Engine? (the question-shape chart)}

\begin{center}
\small
\renewcommand{\arraystretch}{1.25}
\setlength{\tabcolsep}{3pt}
\begin{tabular}{@{}p{0.245\linewidth} p{0.155\linewidth} p{0.305\linewidth} >{\raggedright\arraybackslash}p{0.13\linewidth}@{}}
\toprule
\textbf{Question} & \textbf{Engine} & \textbf{Why it is not a plain tag} & \textbf{Where} \\
\midrule
insert lines, query pointwise \(\min/\max\) & Li Chao & no fixed array; the domain is the tree & \secref{c29-lichao} \\
DP over a whole ancestor path, arbitrary slopes & Li Chao + rollback & the line set changes as you descend & \secref{c29-path} \\
range \(\mathit{chmin}/\mathit{chmax}\) + sum & beats & clamp depends on each value & \secref{c29-beats} \\
range modulo / floor & beats (max-prune) & only elements \(\ge x\) move & \secref{c29-beats} \\
assign/merge runs, random data & Chtholly (ODT) & amortised on generated orderings & \secref{c29-beats} \\
range apply a \emph{product of maps} & function-composition & node stores a transform, merges by compose & \secref{c29-func} \\
range add / assign, composable & \emph{plain lazy} & it \emph{is} a monoid -- go back & \chref{c15} \\
\bottomrule
\end{tabular}
\end{center}

\begin{insightbox}[One axis decides everything: is the update a composable monoid?]
Chapters~19--22 all live on one side of a line: their range writes \emph{compose}
(\((\text{add }a)\) then \((\text{add }b) = (\text{add }a{+}b)\)) so a tag can be
deferred.  Cross the line -- \(\mathit{chmin}\), modulo, a per-value clamp -- and
composition fails, because the effect depends on the value it lands on.  The
three engines here are the three ways back to almost-logarithmic time once the
tag is gone: \emph{change the domain} (Li Chao indexes the \(x\)-axis, not the
array), \emph{change the node state} (beats keeps \(\max/\mathrm{smax}/\mathrm{cnt}\) and
prunes), or \emph{change what a node means} (function-composition stores a map
and merges by composing).  The rest of the chapter is cashing in that one
distinction (\(\secref{c29-boundary}\)'s pick table).
\end{insightbox}

% ---------------- Taxonomy chart ----------------
\begin{figure}[H]
\centering
\begin{tikzpicture}[diagram, scale=0.95, transform shape,
    mnode/.style={circle, draw=#1, line width=0.9pt, fill=#1!8, minimum size=8mm,
                  font=\scriptsize\bfseries, text=#1!60!black}]
  \node[mnode=primaryblue] (a) at (0,0)   {line};
  \node[annot, font=\scriptsize, anchor=west, text width=9.4cm, align=left] at (1.35,0) {a Li Chao container indexes the x-axis; insert a line keeping the mid-winner, recurse the loser to one side; a point query is a root-to-leaf max (23.1)};
  \node[mnode=successgreen] (b) at (0,-1) {path};
  \node[annot, font=\scriptsize, anchor=west, text width=9.4cm, align=left] at (1.35,-1) {the same container rides a DFS: insert the node's line on entry, roll it back on exit, so a query sees exactly its ancestor lines (23.2, COT6)};
  \node[mnode=difforange]   (c) at (0,-2) {beats};
  \node[annot, font=\scriptsize, anchor=west, text width=9.4cm, align=left] at (1.35,-2) {store (max, second-max, count-of-max); a chmin prunes if x>=max, applies at the node if smax<x<max, else descends -- amortised O(log n) blends (23.3)};
  \node[mnode=warningred]   (d) at (0,-3) {odt};
  \node[annot, font=\scriptsize, anchor=west, text width=9.4cm, align=left] at (1.35,-3) {the Chtholly tree keeps maximal equal-runs in an ordered map; assign collapses runs, so random data amortises (23.3, CF 896C)};
  \node[mnode=softpurple]   (e) at (0,-4) {map};
  \node[annot, font=\scriptsize, anchor=west, text width=9.4cm, align=left] at (1.35,-4) {a node stores a transform (affine or a matrix) and merges by composition; a range applies a product of maps -- the general case of Ch.20's lazy (23.4)};
  \node[mnode=deepblue]     (f) at (0,-5) {pick};
  \node[annot, font=\scriptsize, anchor=west, text width=9.4cm, align=left] at (1.35,-5) {composable? a lazy tag (Ch.20). point-extremum over lines? Li Chao. per-value clamp? beats/ODT. a node is a function? compose. dynamic hull? Ch.25 (23.5)};
\end{tikzpicture}
\caption{The chapter's six-state taxonomy.  \emph{line} and \emph{path} are the
Li Chao engine read statically then on a tree; \emph{beats} and \emph{odt} are
the two amortised answers to a non-composable clamp; \emph{map} is the
function-composition generalisation of a lazy tag; \emph{pick} is the boundary
that chooses among them and hands off to the neighbouring chapters.  Every solved
problem sits on one row: CF~631E on \emph{line}, SPOJ~COT6 on \emph{path},
CF~896C on \emph{odt}, CF~750E on \emph{map}, the Summit CF~932F on
\emph{line}+\emph{path} merged.}
\label{fig:c29-taxonomy}
\end{figure}

% ---------------- Who owns what ----------------
\subsection*{Who Owns What (so nothing is learned twice)}

\begin{center}
\small
\renewcommand{\arraystretch}{1.25}
\begin{tabular}{@{}p{0.30\linewidth} p{0.28\linewidth} p{0.36\linewidth}@{}}
\toprule
\textbf{Idea} & \textbf{Owner} & \textbf{This chapter's share} \\
\midrule
build / point update / range fold & Ch.~19 (\chref{c13}) & the engines reuse the skeleton; only the node state or the merge changes \\
lazy tag as a compose monoid & Ch.~20 (\chref{c15}) & 23.4 \emph{is} that monoid generalised (map compose); 23.3 is the escape hatch when it fails -- back-cited, never re-derived \\
static \(k\)-th / count on a path & Ch.~21 (\chref{c19}) & forward-cited rival; the range \(k\)-th here is answered by ODT, not persistence \\
path = \(O(\log n)\) chain ranges & Ch.~22 (\chref{c16}) & contrasted in 23.2: a whole ancestor path queried at ONCE by a container, not tiled \\
modulo-prune beats, CF 438D & Ch.~19 (\chref{c13}) & PROSE re-treat only; 438D is owned there -- never re-linked \\
dynamic convex hull / Li Chao with deletions & Ch.~25 (\chref{c23}) & forward-cited boundary; rollback here is offline-on-DFS, not a fully dynamic hull \\
offline range / path queries & Ch.~27 (\chref{c21}) & forward-cited when order is free; no code here \\
\bottomrule
\end{tabular}
\end{center}

% ---------------- Recognition guide ----------------
\begin{recognitionbox}[Pattern card -- ``a range write that is not a composable tag'']
\reckey{updates of the form ``set \(a_i\leftarrow\min(a_i,x)\)/\(\max(a_i,x)\)'',
``\(a_i\mathrel{\%}=x\)'', ``\(a_i\leftarrow\lfloor a_i/x\rfloor\)'' interleaved
with range \(\sum/\max\) queries (beats); ``insert a line, ask the extreme at a
point \(x\)'' or a DP transition \(f(i)=\min_j\{m_j x + b_j\}\) with arbitrary
slopes (Li Chao); ``apply a whole sequence of maps over a range'' where the maps
compose (function-composition).  If the update is add/assign/multiply/flip only,
it \emph{is} a monoid -- use Chapter~20 instead.}
\recuse{
  \item pointwise extremum over a family of lines, arbitrary slopes -- Li Chao (\(\secref{c29-lichao}\)); on a tree path -- rollback (\(\secref{c29-path}\));
  \item range \(\mathit{chmin}/\mathit{chmax}/\)modulo with a running sum -- beats (\(\secref{c29-beats}\)); heavy \emph{assign} on random data -- ODT (\(\secref{c29-beats}\));
  \item a range is rewritten by composing per-element transforms -- function-composition (\(\secref{c29-func}\)).
}
\recnotuse{
  \item the update is a single composable op (add, assign, multiply, flip) -- Ch.~20's lazy tag (\chref{c15});
  \item the update is a plain point write and the query a fold -- Ch.~19's tree (\chref{c13});
  \item the convex hull must support \emph{deletions} / arbitrary insertion online -- link-cut / dynamic hull, Ch.~25 (\chref{c23});
  \item the queries are offline and the order is free -- Mo, Ch.~27 (\chref{c21}).
}
One-line test: \emph{the update's effect depends on the value it lands on (clamp,
modulo) or there is no fixed array (a moving line envelope) -- then this chapter,
not a lazy tag.}
\end{recognitionbox}

% ---------------- Constraint -> engine ----------------
\subsection*{Reading the Constraints (which engine the limits want)}

\begin{center}
\small
\renewcommand{\arraystretch}{1.25}
\begin{tabular}{@{}p{0.32\linewidth} p{0.40\linewidth} p{0.22\linewidth}@{}}
\toprule
\textbf{Constraint signature} & \textbf{What it rules in/out} & \textbf{Engine} \\
\midrule
query extreme of lines at integer \(x\in[1,C]\), \(C\le10^{9}\) & no array; index the domain, not positions & \secref{c29-lichao} \\
DP \(f(u)=\min\) over ancestors of a line at a node-dependent \(x\) & a container per root-path with rollback & \secref{c29-path} \\
\(N,Q\le10^{5}\); range chmin/chmax/add/sum mixed & a tag can't compose a clamp; the node needs \((\max,\mathrm{smax})\) & \secref{c29-beats} \\
ops generated at random; range assign + \(k\)-th + power-sum & worst-case loose, random runs collapse & \secref{c29-beats} \\
map a range by \(x\mapsto ax+b\) or a matrix, product queries & node = a transform, merge = compose & \secref{c29-func} \\
\bottomrule
\end{tabular}
\end{center}

% ---------------- Notation ----------------
\subsection*{Notation Used Only Here}

The array is \(A=(a_1,\dots,a_n)\); ranges \([l,r]\) are index intervals; a
\emph{segment-tree node} covers an interval \(I\) and \(\mathit{seg}(I)\) is its
stored aggregate.  For \emph{beats} (\(\secref{c29-beats}\)):
\(\max(I)\) and \(\mathrm{smax}(I)\) are the largest and second-largest value in \(I\),
\(\mathrm{cnt}(I)\) the multiplicity of \(\max(I)\); \(\mathit{chmin}(x)\) sends every
value to \(\min(\text{value},x)\).  For \emph{Li Chao} (\(\secref{c29-lichao}\)):
a line is \(\ell(x)=m x + b\), the domain is \([1,C]\) (integer \(x\)), and
\(\ell_1\!\le\!\ell_2\) at \(x\) means \(\ell_1(x)\le\ell_2(x)\).  For
\emph{function-composition} (\(\secref{c29-func}\)): a transform is written
\(f\!\circ\!g\) for ``apply \(g\) then \(f\)'', an affine map is
\(x\mapsto ax+b\), and matrices multiply on the appropriate side (a
\((\min,+)\)-product for the automaton tree).  \(\bigO{\log C}\) is over the
\emph{domain} size, \(\bigO{\log n}\) over the array length.  The lazy
composition monoid and the \code{push} discipline are Chapter~20's
(\(\secref{c15-compose}\)), reused verbatim -- never re-derived.  Every section runs on
\(A\) above; if a table here and a trace there disagree, this opener is
authoritative.
